\documentclass[11pt]{article}

\usepackage[margin=1in]{geometry}
\usepackage{amsmath}
\usepackage{amssymb}
\usepackage{booktabs}
\usepackage{longtable}
\usepackage[T1]{fontenc}
\usepackage{microtype}

\newcommand{\SDU}{\mathrm{SDU}}
\newcommand{\NaN}{\mathrm{NaN}}
\newcommand{\model}{\mathcal{M}}

\title{LuSEE Spectrometer Gain Conversion}
\author{}
\date{}

\begin{document}
\maketitle

\section{Indices and stored quantity}

Let $t$ index spectra, $p\in\{0,\ldots,15\}$ index spectrometer
products, $c\in\{0,\ldots,3\}$ index ADC channels, and $f$ denote a
frequency bin.  Let $C_{tpf}$ be the decoded firmware value and let
$b^{\mathrm{actual}}_{tp}$ be the corresponding realized bit-slice value.
Ingestion always restores the decoded normal spectrum to the bit-31
accumulator convention:
\begin{equation}
  X^{\SDU}_{tpf}
  = \operatorname{ldexp}\!\left(C_{tpf},
      b^{\mathrm{actual}}_{tp}-31\right)
  = C_{tpf}\,2^{b^{\mathrm{actual}}_{tp}-31}.
  \label{eq:bitslice}
\end{equation}
Only $X^{\SDU}$ is stored in \texttt{/spectra/data}.  The realized
bit-slice metadata is retained as an $(N,16)$ array.  This operation is
not optional and the requested bit-slice value is never substituted for a
missing realized value.

\section{Vendored coefficient table}

The concrete fit artifact is \texttt{lusee/data/gain/alpha\_refit.csv}.  It
has 318 data rows and exactly the columns
\texttt{gain\_setting}, \texttt{component}, \texttt{model}, \texttt{term},
and \texttt{alpha\_refit}.  Its row counts are:
\begin{center}
\begin{tabular}{lrrr}
\toprule
family & PC1 & PC2 & total \\
\midrule
linear    & 64 & 68 & 132 \\
quadratic & 91 & 95 & 186 \\
\midrule
total     & 155 & 163 & 318 \\
\bottomrule
\end{tabular}
\end{center}
There are 12 gain/channel settings, $L0$ through $H3$, and both fitted
families are present for each of PC1 and PC2.  The table is sparse: its term
vocabulary is the intercept, five linear telemetry variables, and the three
temperature products shown below (with ADC0 or ADC1 selected by channel).
The CSV itself, rather than a presumed polynomial template, is the definitive
list of nonzero coefficients.

For completeness, Table~\ref{tab:sparse-masks} gives the concrete sparse
term mask of every fitted group.  The abbreviations are
\begin{equation*}
  \begin{aligned}
    0&=1, & F&=T_{\mathrm{FPGA}}, & A&=T_{A,c}, \\
    V&=V_{1.8}, & D&=V_{1.2}, & I&=I_{1.8}, \\
    F^2&=T_{\mathrm{FPGA}}^2, & A^2&=T_{A,c}^2,
      & FA&=T_{\mathrm{FPGA}}T_{A,c}.
  \end{aligned}
\end{equation*}
Each listed term has its own CSV value
$\lambda^{(m,r)}_{sck}$; omitted terms have coefficient zero.

\begin{longtable}{@{}llp{0.32\textwidth}p{0.40\textwidth}@{}}
\caption{Exact nonzero term masks in \texttt{alpha\_refit.csv}.}
\label{tab:sparse-masks}\\
\toprule
setting & PC & linear family & quadratic family \\
\midrule
\endfirsthead
\toprule
setting & PC & linear family & quadratic family \\
\midrule
\endhead
L0 & 1 & $0,F,A,D,I$ & $F,A,D,F^2,A^2,FA$ \\
L0 & 2 & $0,F,A,V,D,I$ & $0,F,A,V,D,I,F^2,A^2,FA$ \\
L1 & 1 & $0,F,A,V,D,I$ & $0,F,A,V,D,I,F^2,A^2,FA$ \\
L1 & 2 & $0,F,A,V,D,I$ & $0,F,A,D,I,F^2,A^2,FA$ \\
L2 & 1 & $0,F,A,D,I$ & $F,A,V,D,I,F^2,A^2,FA$ \\
L2 & 2 & $0,F,A,V,D,I$ & $0,F,A,V,D,I,F^2,A^2,FA$ \\
L3 & 1 & $0,F,A,V,D,I$ & $0,F,A,V,D,I,F^2,A^2,FA$ \\
L3 & 2 & $F,A,V,D,I$ & $F,D,I$ \\
M0 & 1 & $0,F,A,V,D,I$ & $F,A,I,F^2,A^2,FA$ \\
M0 & 2 & $0,F,A,V,D,I$ & $0,F,A,V,D,I,F^2$ \\
M1 & 1 & $0,A,V,D,I$ & $0,F,A,V,D,I,F^2,A^2,FA$ \\
M1 & 2 & $0,F,A,V,D,I$ & $0,F,A,V,D,I,F^2,A^2,FA$ \\
M2 & 1 & $0,F,A,V,D,I$ & $0,F,A,V,D,I,F^2,A^2,FA$ \\
M2 & 2 & $F,A,V,D,I$ & $0,F,A,V,D,I,F^2,A^2,FA$ \\
M3 & 1 & $0,A,V,D,I$ & $0,A,I,F^2$ \\
M3 & 2 & $0,F,A,V,D,I$ & $0,F,A,V,D,I,F^2,A^2,FA$ \\
H0 & 1 & $0,F,A,D,I$ & $F,A,D,I,F^2,A^2,FA$ \\
H0 & 2 & $0,F,A,V,D$ & $0,F,A,V,D,F^2,A^2,FA$ \\
H1 & 1 & $0,F,A,V,D,I$ & $F,A,V,D,F^2,A^2,FA$ \\
H1 & 2 & $0,F,A,V,I$ & $0,F,A,V,D,I,F^2,A^2,FA$ \\
H2 & 1 & $F,A,V,D$ & $0,F,A,V,D,I,F^2,A^2,FA$ \\
H2 & 2 & $0,F,A,V,D,I$ & $F,A,D,I,F^2,A^2,FA$ \\
H3 & 1 & $0,F,A,D,I$ & $0,F,A,V,D,I,F^2,FA$ \\
H3 & 2 & $0,F,A,V,D,I$ & $0,F,A,V,I,F^2,A^2,FA$ \\
\bottomrule
\end{longtable}

\section{Telemetry regression}

For channel $c$, define
\begin{equation}
  T_{A,c} =
  \begin{cases}
    T_{\mathrm{ADC0}}, & c\in\{0,1\},\\
    T_{\mathrm{ADC1}}, & c\in\{2,3\},
  \end{cases}
\end{equation}
and the common feature vector
\begin{equation}
  \boldsymbol{\phi}_c(\mathbf z_t) =
  \begin{pmatrix}
    1 & T_{\mathrm{FPGA}} & T_{A,c} & V_{1.8} & V_{1.2} & I_{1.8}
      & T_{\mathrm{FPGA}}^2 & T_{A,c}^2
      & T_{\mathrm{FPGA}}T_{A,c}
  \end{pmatrix}^{\!\mathsf T}_{t}.
  \label{eq:features}
\end{equation}
The coefficient table \texttt{alpha\_refit.csv} contains a sparse subset
of these terms for gain setting $s\in\{L,M,H\}$, channel $c$, principal
component $r\in\{1,2\}$, and fitted family
$m\in\{\mathrm{linear},\mathrm{quadratic}\}$.  Define a missing CSV row to
have coefficient zero.  Then the predicted score is
\begin{equation}
  q^{(m_r)}_{tscr}
  = \sum_k \lambda^{(m_r,r)}_{sck}\,
      \phi_{c,k}(\mathbf z_t).
  \label{eq:pc-score}
\end{equation}
For example, the FPGA-temperature contribution is represented symbolically by
\begin{equation*}
  \lambda^{(m_r,r)}_{sc,\mathrm{THERM\_FPGA}}\,
  T_{\mathrm{FPGA},t}.
\end{equation*}
No fitted numerical coefficient is part of this definition.

The labels ``linear'' and ``quadratic'' select rows of the fitted artifact.
They do not dictate which terms must be present.  In particular, a row group
labelled quadratic may retain linear terms, quadratic terms, both, or only a
screened subset.  PC1 and PC2 select their families independently:
\begin{equation}
  (m_1,m_2) \in
  \{\mathrm{linear},\mathrm{quadratic}\}^2.
\end{equation}

\section{Gain reconstruction}

At gain setting $s$, channel $c$, and PCA anchor-frequency index $a$, the
gain-model artifact supplies a mean $\mu_{sca}$ and two eigenvectors
$v_{sc1a}$ and $v_{sc2a}$.  The reconstructed gain is
\begin{equation}
  G_{tca}
  = \mu_{sca}
    +q^{(m_1)}_{tsc1}v_{sc1a}
    +q^{(m_2)}_{tsc2}v_{sc2a}.
  \label{eq:gain-anchor}
\end{equation}
A cubic spline, with extrapolation disabled, evaluates this quantity on the
spectrometer frequency grid:
\begin{equation}
  G_{tc}(f)
  = \mathcal S_{sc}\!\left[
      \mu_{sc\cdot}
      +q^{(m_1)}_{tsc1}v_{sc1\cdot}
      +q^{(m_2)}_{tsc2}v_{sc2\cdot}
    \right](f).
  \label{eq:gain-interp}
\end{equation}

\section{Conversion to physical spectral density}

For an auto product $p=c$, the input-referred power spectral density in the
native gain-model units is
\begin{equation}
  P^{\mathrm{nV}^2/\mathrm{Hz}}_{tcf}
  = \frac{X^{\SDU}_{tcf}}{G_{tc}(f)}.
  \label{eq:auto-psd}
\end{equation}
For either the real or imaginary product of a cross between channels $c$ and
$d$, it is
\begin{equation}
  P^{\mathrm{nV}^2/\mathrm{Hz}}_{t,(cd),f}
  = \frac{X^{\SDU}_{t,(cd),f}}
         {\sqrt{G_{tc}(f)G_{td}(f)}}.
  \label{eq:cross-psd}
\end{equation}
The SI-unit result is
\begin{equation}
  P^{\mathrm{V}^2/\mathrm{Hz}} = 10^{-18}
  P^{\mathrm{nV}^2/\mathrm{Hz}}.
\end{equation}

For compatibility with the production reference, the amplitude-spectral-
density view is
\begin{align}
  A^{\mathrm{nV}/\sqrt{\mathrm{Hz}}}_{tcf}
    &= \sqrt{X^{\SDU}_{tcf}/G_{tc}(f)},
      &&\text{auto}, \label{eq:auto-asd}\\
  A^{\mathrm{nV}/\sqrt{\mathrm{Hz}}}_{t,(cd),f}
    &= \operatorname{sgn}(X^{\SDU}_{t,(cd),f})
       \sqrt{\frac{|X^{\SDU}_{t,(cd),f}|}
                    {\sqrt{G_{tc}(f)G_{td}(f)}}},
      &&\text{cross component}. \label{eq:cross-asd}
\end{align}
Equation~\eqref{eq:cross-asd} is a signed display convention for separate
real and imaginary components; equation~\eqref{eq:cross-psd} is the
physically linear cross-power representation.

Non-finite telemetry or data, a non-positive reconstructed gain, a negative
auto power, or a frequency outside the anchor range produces $\NaN$.  The
conversion is evaluated lazily when requested and is not written back to the
HDF5 file.  Returned arrays are decorated with either
\texttt{nV/sqrt(Hz)}, \texttt{nV\^{}2/Hz}, or \texttt{V\^{}2/Hz}; stored
normal spectra are decorated with \texttt{SDU}.

\section{Frozen implementation contract}

The implementation intentionally freezes the vendored PCA arrays and the
rows of \texttt{alpha\_refit.csv}.  A process-wide selection chooses
$(m_1,m_2)$ once per public conversion call.  Dense coefficient and feature
arrays are used for the polynomial evaluation, while reconstructed gains are
vectorized over time, channel, and frequency.  This document specifies
reproduction of that artifact; it does not independently resolve the
pre-amplifier reference-temperature or FFT equivalent-noise-bandwidth
conventions used when the artifact was produced.

\end{document}
